%% Introduction %%

Bien que ces deux deux services suivent les mêmes missions de conservation et communication, ils s'opposent à plusieurs niveaux. Là où le dépôt légal applique ces missions aux documents captés ou déposés au titre du dépôt légal, le service des archives professionnelles concerne la conservation et l’exploitation commerciale des archives de l’ORTF et les archives confiées par la loi à l’INA sur lesquels l’institution a un droit d’exploitation. Au delà de leurs fonds, ces services s'opposent aux archives du dépôt légal dans leur organisation. Les archives professionnelles fonctionnent en « stock », elles gèrent des stocks, des programmes, des émissions à l’unité, qu’elles exploitent, au contraire du dépôt légal qui a une gestion en « flux », selon les tranches horaires, les grilles de programmes\footcite[][23]{varra2023}. Les archives professionnelles conservent un programme dans la meilleures qualité possible, le dépôt légal le conserve autant de fois qu'il a été diffusé. C'est la manifestation des publics et usages différents qu'ils servent et que nous avons déjà évoqués plus tôt. Les archives professionnelles communiquent avec les professionnels de l’audiovisuel : les chaînes de télévisions, les journaux télévisés, ils ont donc besoin d'accéder à leur contenu par sujet, par programme unique. Le dépôt légal communique avec les chercheurs de l'INAthèque, qui peuvent aussi bien travailler sur un programme que sur une chaîne ou sur une tranche dans leur ensemble. Pour répondre à ces usages différents, les deux services adoptent deux politiques documentaires différentes. 

\subsection{Des données de descriptions divergentes}

Aujourd’hui, ces deux services ont disparu à la suite de plusieurs réorganisations des services, dont la dernière, en 2020, qui répartit ces deux services dans la « Direction des Patrimoines »\footnote{(Voir organigrammes de la direction des Patrimoines :~\ref{ann:org},\nameref{ann:org}), p.~\pageref{ann:org}.}. Cependant, ils transparaissent toujours de diverses manières au sein de la politique documentaire en vigueur. Nous avons déjà vu que leurs objectifs et leur public n’était pas les mêmes, mais nous n’avons pas vu de quelle façon cela se répercutait sur le traitement documentaire. Pour illustrer nos propos, prenons pour exemple deux notices d’une même émission\footnote{(Pour consulter les notices voir :~\ref{ann:notices},\nameref{ann:notices}), p.~\pageref{ann:notices}.}, une notice provient de Totem, l'outil de consultation de la base de données des archives professionnelles, l’autre vient de DLTV, la base de données du dépôt légal. Évidemment, les deux notices présentent des informations de description identiques telles que le titre, la date de diffusion, le genre d’émission, le générique et même le sommaire. Mais les informations sur la diffusion diffèrent. Du côté professionnel, nous trouvons la date, l’heure mais l’information importante du côté dépôt légal est que cette émission est en première diffusion. Cette information a peu de valeur dans le cadre d’une commercialisation mais elle en a une à la fois pour les TGDM chargés de déterminer le type de traitement à appliquer à cette émission et pour les chercheurs de l’INAthèque. Les autres informations de diffusion sont présentes dans la fiche Hyperbase mais elles sont réparties entre différents champs afin de permettre la recherche par ce champ. Les autres données de description le confirment. Dans la fiche des archives professionnelles, nous retrouvons les informations concernant le statut de numérisation, essentiel pour partager ce document, la dévolution, sur laquelle nous allons revenir, et surtout des informations précises sur les fichiers, pour la diffusion de ce document. Toutes ces informations sont absentes de la fiche du dépôt légal. Au contraire, nous trouvons dans la fiche du dépôt légal des informations sur les documents d’accompagnements, pouvant compléter une recherche et les audiences permettant, entre autres, de réaliser des statistiques. Usage courant pour un chercheur mais improbable pour un professionnel de l'audiovisuel.

\subsection{Un fossé entre les services de conservation} 

En comparant ces deux fiches, il est évident que même actuellement, des années après la réorganisation des services, deux politiques documentaires cohabitent. Pourtant, l’écart entre les deux services patrimoniaux est bien moins important qu’il ne le fut. À la création du dépôt légal, il est nécessaire de trouver de nouveaux locaux pour accueillir ce nouveau service. Le service du dépôt légal investit des locaux à une centaine de mètres des locaux historiques de l’INA à Bry-sur-Marne\footcite[][24]{varra2023}, creusant un fossé entre les deux services, même physiquement. C’est un ressenti évoqué par une documentaliste interrogée, anciennement au dépôt légal, qui souligne l’écart entre les deux services. L’accès aux notices, aux images et aux sons des documents du service des archives professionnelles étaient compliquées avant leur numérisation car ils n’avaient même pas accès à l'outil de consultation du service\footnote{(Voir entretien d'une documentaliste de l'INA sur la différence entre le dépôt légal et les archives professionnelles~\ref{ann:transcriptions},\nameref{ann:transcriptions}), p.~\pageref{ann:transcriptions}.}. Dans ces circonstances, on peut difficilement imaginer une bonne communication sur les politiques documentaires. D'autant plus que ce n'était pas la première fois de l'histoire de la conservation des documents audiovisuels que différents services de l’INA furent séparés géographiquement. Avant même sa création, la cinémathèque était séparée entre les services « Film d’Actualités » à Cognacq-Jay et « Film Production » aux Buttes Chaumont. Ils seront réunis pour une courte durée à la création de l’INA, à l'occasion de laquelle ils réunissent leur fiches papiers, non sans difficulté\footcite[]{saintville1994}, avant de se séparer à nouveau puis d’être réunis à Bry-sur-Marne. Nous pouvons donc imaginer, en conséquence, les nouvelles incohérences générées dans les politiques documentaires. Une fois de plus, nous savons qu'elles existent, mais difficile de pointer avec exactitude lesquelles sont affectées tant la communication étaient compromises et l'archivage des pratiques documentaires absentes. 

Nous pouvons terminer en soulignant toutes les erreurs d'origine humaines : fautes de frappes, erreurs d'orthographes notamment des noms, omission volontaire ou non. Le contexte de description de ces notices peut également avoir un impact. Dans le cas des prises d'antenne, si une grève a lieu, il n'y a pas de description et d'indexation. Certaines notices sont indiquées \enquote{non indéxé parce qu'en grève}\footnote{(Voir entretien d'une documentaliste sur l'impact de la politique documentaire sur les données de description~\ref{ann:transcriptions},\nameref{ann:transcriptions}), p.~\pageref{ann:transcriptions}.}. Cela dit, le manque de communication n’est pas ignoré, à plusieurs reprises, il sera tenté d'harmoniser les méthodes. Des référentiels, des thésaurus sont mis en place nous pouvons plus particulièrement souligner la création de la DDCOLL, direction déléguée des collections, en 2010 qui met en place une structure transversale « méthodes et qualités » et qui donne lieu à la réunification des services, évoqués plus haut\footcite[][29]{varra2023}, ainsi que la création du « lac de données » sur lequel nous reviendrons. 

%% Conclusion %%

Ces perturbations d’ordre institutionnel ne trouvent pas leur pareil au sein des institutions patrimoniales ayant pour mission le dépôt légal. Il n’y ni au sein du CNC, ni au sein de la BNF de distinction entre les documents du dépôt légal et les documents commercialisés. C’est une exception que l’on doit au statut d’EPIC (Établissement public à caractère industriel et commercial en France) de l’INA que l'on ne retrouve pas à la BnF ou au CNC. Cette exception est responsable de cet écart actuel entre le traitement documentaire réalisé sur des archives professionnelles et celui sur les archives du dépôt légal. C'est aujourd'hui un enjeu d'accessibilité car il peut-être difficile de savoir quelles informations peut-on trouver et où ? 